% =====================================================================
%  Pakete und Formatvorgaben laut FH-Burgenland-Vorlage (Seminararbeit)
%  (abgeleitet aus dem MSc-Thesis-Template, 00_A_preamble/thesis.tex)
% =====================================================================
\usepackage{scrhack}
\usepackage{etoolbox}                       % \ifdefempty fuer optionale Deckblatt-Zeilen

% -- Seitenlayout: A4, links/rechts 2,5 cm, oben 2,5 cm, unten 2 cm, Bundsteg 1 cm
\usepackage[left=2.5cm,right=2.5cm,top=2.5cm,bottom=2cm,bindingoffset=1cm,headsep=1.25cm,footskip=1.25cm]{geometry}

% -- Kopf-/Fusszeile: nur Seitenzahl rechts unten (10 pt), keine Kopfzeile
\usepackage[automark]{scrlayer-scrpage}
\clearpairofpagestyles
\ofoot*{\fontsize{10pt}{12pt}\selectfont\pagemark}
\pagestyle{scrheadings}

% -- Sprache, Zeichensatz, Schrift
\usepackage[ngerman]{babel}
\usepackage[T1]{fontenc}
\usepackage{mathpazo}                       % "Book Antiqua" entspricht Palatino
\usepackage[final]{microtype}
\usepackage{csquotes}                       % benoetigt fuer babel + biblatex-apa
\usepackage{xurl}

% -- Grafiken, Tabellen, Mathematik, Code
\usepackage{graphicx}
\usepackage[table,xcdraw]{xcolor}
\usepackage{booktabs}
\usepackage{tabularx}
\usepackage{multirow}
\usepackage{array}
\usepackage{amsmath}
\usepackage{listings}
\usepackage{algorithm2e}
\usepackage{verbatim}
\usepackage{fmtcount}
\usepackage{changepage}
\usepackage{pifont}
\newcommand{\cmark}{\ding{51}}
\newcommand{\xmark}{\ding{55}}
\newcommand{\omark}{\ding{72}}
\usepackage{lipsum}                         % Blindtext - kann entfernt werden

% -- Abkuerzungen / Glossar
\usepackage[printonlyused, nohyperlinks]{acronym}
\usepackage[record=only,acronym,nonumberlist,nopostdot=false,toc=true,style=super]{glossaries-extra}

% -- Abbildungen, Tabellen und Formeln fortlaufend nummerieren (wie in der Word-Vorlage:
%    "Abbildung 1", "Tabelle 1", "Formel 1" statt "Abbildung 2.1")
\counterwithout{figure}{chapter}
\counterwithout{table}{chapter}
\counterwithout{equation}{chapter}

% -- Beschriftungen: 10 pt, linksbuendig
\usepackage[justification=RaggedRight,singlelinecheck=false,font=footnotesize]{caption}

% -- Gliederung: max. 3 Ebenen im Inhaltsverzeichnis (z.B. 3.3.1)
\setcounter{secnumdepth}{3}
\setcounter{tocdepth}{2}

% -- Schriftgroessen der Ueberschriften laut Vorlage
\setkomafont{disposition}{\bfseries}
\addtokomafont{title}{\fontsize{24}{24}\selectfont}         % Titel der Arbeit
\addtokomafont{subtitle}{\fontsize{14}{16}\selectfont}      % uebriger Deckblatt-Text
\addtokomafont{chapter}{\fontsize{16}{16}\selectfont}       % Ebene 1
\addtokomafont{section}{\fontsize{14}{16}\selectfont}       % Ebene 2
\addtokomafont{subsection}{\fontsize{14}{16}\selectfont}    % Ebene 3
\addtokomafont{subsubsection}{\fontsize{13}{16}\selectfont} % Ebene 4

% -- Tabellen: 10 pt (erste Zeile 11 pt -> {\fontsize{11pt}{12pt}\selectfont ...})
\let\oldtabular\tabular
\renewcommand{\tabular}{\fontsize{10pt}{14pt}\selectfont\oldtabular}

% -- Absaetze: 12 pt, Zeilenabstand 16 pt, 6 pt Abstand, Blocksatz, kein Einzug
\linespread{1.1}                            % 12 pt Palatino -> ca. 16 pt Zeilenabstand
\setlength{\parskip}{6pt}
\setlength{\parindent}{0pt}
\renewcommand*{\chapterheadstartvskip}{\vspace*{0cm}}   % Kapitel beginnt oben auf der Seite

\RedeclareSectionCommand[beforeskip=0pt,afterskip=.5\baselineskip plus .1\baselineskip minus .167\baselineskip]{chapter}
\RedeclareSectionCommand[beforeskip=-1\baselineskip,afterskip=.2\baselineskip]{section}
\RedeclareSectionCommand[beforeskip=-.75\baselineskip,afterskip=.5\baselineskip]{subsection}
\RedeclareSectionCommand[beforeskip=-.75\baselineskip,afterskip=.25\baselineskip]{subsubsection}
\RedeclareSectionCommand[beforeskip=.5\baselineskip,afterskip=-1em]{paragraph}
\RedeclareSectionCommand[beforeskip=-.5\baselineskip,afterskip=-1em]{subparagraph}

% -- Quellcode-Darstellung
\definecolor{codegreen}{rgb}{0,0.6,0}
\definecolor{codegray}{rgb}{0.5,0.5,0.5}
\definecolor{codepurple}{rgb}{0.58,0,0.82}
\definecolor{backcolour}{rgb}{0.95,0.95,0.92}
\lstdefinestyle{mystyle}{
    backgroundcolor=\color{backcolour},
    commentstyle=\color{codegreen},
    keywordstyle=\color{magenta},
    numberstyle=\tiny\color{codegray},
    stringstyle=\color{codepurple},
    basicstyle=\ttfamily\footnotesize,
    breakatwhitespace=false,
    breaklines=true,
    captionpos=b,
    keepspaces=true,
    numbers=left,
    numbersep=5pt,
    showspaces=false,
    showstringspaces=false,
    showtabs=false,
    tabsize=2
}
\lstset{style=mystyle}

% -- Literatur: APA (Autor, Jahr, Seite) mit Biber
\usepackage[
    style=apa,
    backend=biber,
    sorting=nyt,
    language=german,
    sortlocale=de-DE,
    natbib=true,
    uniquename=false,
    uniquelist=false,
    maxbibnames=1000,
    urldate=long,
]{biblatex}
\addbibresource{references.bib}

\usepackage{hyperref}
